\begin{center}
  {\LARGE\bfseries STCET ECE Secures NBA Accreditation\par}
  \vspace{0.45em}
  {\large\bfseries\color{IEEEblue}
  Recognition for the 2026--2028 cycle highlights outcome-based learning, practical education and continuous academic improvement.\par}
  \vspace{0.5em}
  {\small\itshape By the Editorial Desk}
\end{center}

\begin{tcolorbox}[
  colback=SoftBlue,
  colframe=IEEEblue!75!black,
  boxrule=0.6pt,
  arc=0mm,
  left=3mm,right=3mm,top=2mm,bottom=2mm
]
  \small\RaggedRight
  \textbf{At a glance:} The National Board of Accreditation has accredited the B.Tech. programme in Electronics and Communication Engineering at St. Thomas' College of Engineering \& Technology for the 2026--2028 cycle. The college's official records also list accreditation cycles for 2019--2022 and 2022--2025.
\end{tcolorbox}

{\small
\noindent
\textbf{Kolkata, August 2026---}The B.Tech. programme in Electronics and Communication Engineering (ECE) at St. Thomas' College of Engineering \& Technology (STCET) has secured accreditation from the National Board of Accreditation (NBA) for the 2026--2028 cycle, providing an independent endorsement of the programme's academic processes, outcome-based teaching and commitment to preparing graduates for a changing engineering profession.

\medskip
\noindent
The recognition follows a detailed assessment of how the programme defines its educational goals, delivers its curriculum, supports practical learning, measures student achievement and uses evidence to improve performance. NBA accreditation is awarded to a specific academic programme rather than automatically to an entire institution, making the review directly relevant to the experience and outcomes of ECE students.

\medskip
\noindent
The evaluation extends beyond examination results or a one-time inspection of facilities. It considers programme and course outcomes, faculty strength, student performance, laboratories, computing and library resources, academic support, governance and continuous improvement. At the centre of the process is a practical question: what can graduates understand, analyse, design, build, communicate and improve by the time they complete the programme?

\medskip
\noindent
That approach reflects Outcome-Based Education (OBE), in which teaching and assessment are planned around the capabilities expected of graduates. For ECE students, these capabilities include applying mathematics and engineering science, analysing electronic and communication systems, conducting experiments, using modern tools, designing within realistic constraints, working in teams and recognising professional, ethical, environmental and societal responsibilities.

\medskip
\noindent
Faculty members are expected to map individual course outcomes to broader graduate attributes and measure attainment through examinations, laboratory work, projects, presentations and other evidence. When the results reveal a gap, the programme must show how it responded. Improvements may include revising experiments, strengthening project reviews, introducing academic support, upgrading equipment or software, and expanding interaction with industry and alumni.

\medskip
\noindent
STCET's ECE programme combines classroom foundations with practical work in electronic devices, digital and analogue circuits, microprocessors, signal processing, embedded systems, VLSI and wireless communication. Laboratory sessions and projects allow students to move from equations and circuit diagrams to measurement, implementation, troubleshooting and complete system design---skills that remain central to both core engineering roles and technology-driven careers.

\noindent
Accreditation preparation is a collective exercise. Faculty members and administrators document curriculum delivery and attainment; technical staff maintain laboratories and equipment; students provide evidence through academic work and feedback; and alumni and employers help the department understand how well graduate abilities meet professional expectations. The process is intended to connect institutional claims with verifiable results.

\medskip
\noindent
The programme is associated with NBA's Tier-II framework, which is used for programmes offered by affiliated institutions. India is a permanent signatory to the Washington Accord through the NBA, placing the country's outcome-based accreditation system within an internationally benchmarked engineering-education movement. However, NBA's official rules restrict Washington Accord mutual-recognition eligibility to programmes accredited under Tier-I. Tier-II accreditation should therefore not be presented as automatic Washington Accord equivalence or as a guarantee of recognition overseas.

\medskip
\noindent
Its wider relevance lies in the use of transparent learning outcomes, evidence-based assessment and continuous improvement---principles widely followed in modern engineering education. Graduates applying for higher studies, professional recognition or employment abroad must still satisfy the requirements of the receiving university, employer, licensing body or country.

\noindent
For students and parents, accreditation offers added assurance that the programme is monitored against defined academic standards. It encourages consistency in course delivery, laboratory practice, assessment and student support while requiring the department to demonstrate that feedback and attainment data lead to real changes.

\medskip
\noindent
The potential benefits are also visible in career preparation. Employers increasingly seek graduates who can combine strong fundamentals with experimentation, design, teamwork and communication. Laboratory exercises, mini-projects, internships, seminars and major projects give students opportunities to demonstrate those abilities and build stronger portfolios for placement interviews.

\medskip
\noindent
ECE graduates now enter fields ranging from semiconductor design, VLSI and embedded systems to telecommunications, signal processing, automation, IoT, software and data-driven technology services. Accreditation can help a department keep its teaching connected to these changing expectations through stakeholder feedback, curriculum enrichment, industry-led sessions and project mentoring. It does not guarantee a job offer or salary, but it supports the academic conditions from which better placement readiness can develop.

\medskip
\noindent
The accreditation can similarly add credibility when graduates apply for higher education, scholarships or research opportunities. Individual outcomes will continue to depend on grades, entrance tests, recommendations, research experience, language requirements and the policies of each university. The recognition is best understood as a quality indicator for the programme, not an automatic passport to admission or employment.

\medskip
\noindent
For the department, the award is not the end of the exercise. Accreditation is time-bound, and the programme must continue collecting evidence, analysing outcomes and addressing weaknesses. Faculty development, laboratory modernisation, industry and alumni participation, research-led teaching and support for students with different learning needs will remain important throughout the cycle.

\medskip
\noindent
The 2026--2028 recognition is therefore both a milestone and a mandate for STCET's ECE programme. Its lasting value will be measured not only by the accreditation certificate, but by the confidence, professional judgement and technical capability of the engineers who graduate from the department.
\par}

\vfill
\noindent
{\footnotesize\sloppy\textit{Source note: Accreditation-cycle information is based on the  STCET NBA accreditation page. Accreditation definitions, OBE principles and the Tier-I qualification for Washington Accord mutual recognition follow the NBA General Manual for Accreditation. Programme details are based on official STCET ECE pages.}\par}